\subsection*{I.4.1 Group-mean normalisation}

Let
\[
\bar h=\frac{1}{K}\sum_{i=1}^K h_i.
\]

\paragraph{(i)}
Since $\mathbb{E}[r_i]=\mathbb{E}[\bar r]=\mu$,
\[
\mathbb{E}[\hat A_i]
=
\frac{\mathbb{E}[r_i-\bar r]}{\sigma}
=0.
\]
Because $X_i=h_i\hat A_i$ and $h_i$ is fixed, linearity of expectation gives
\[
\mathbb{E}[X_i]
=
h_i\,\mathbb{E}[\hat A_i]
=
h_i\cdot 0
=0,
\qquad
\boxed{\mathbb{E}[X_i]=0}.
\]
Summing over the group and again using linearity,
\[
\mathbb{E}[\hat g]
=
\mathbb{E}\!\left[\frac{1}{K}\sum_{i=1}^K X_i\right]
=
\frac{1}{K}\sum_{i=1}^K\mathbb{E}[X_i]
=
\frac{1}{K}\sum_{i=1}^K 0
=0,
\qquad
\boxed{\mathbb{E}[\hat g]=0}.
\]
Every expectation in this question is taken conditional on the sampled
responses $a_1,\dots,a_K$, which are held fixed; that is, $\mathbb{E}[\,\cdot\,]$
abbreviates $\mathbb{E}[\,\cdot\mid a_1,\dots,a_K]$, so each score $h_i$ is
constant and only the artificial reward noise is random. In the usual
policy-gradient problem the responses are sampled from the policy and their
rewards depend on the sampled responses, so the policy gradient need not vanish.

\paragraph{(ii)}
The common random term is the sample mean $\bar r$, appearing in every
\[
X_i=\frac{h_i}{\sigma}(r_i-\bar r).
\]
For $i\neq j$,
\begin{align*}
\operatorname{Cov}(r_i-\bar r,r_j-\bar r)
&=
\operatorname{Cov}(r_i,r_j)
-\operatorname{Cov}(r_i,\bar r)
-\operatorname{Cov}(\bar r,r_j)
+\operatorname{Var}(\bar r)\\
&=
0-\frac{\sigma^2}{K}-\frac{\sigma^2}{K}
+\frac{\sigma^2}{K}\\
&=
-\frac{\sigma^2}{K}.
\end{align*}
Hence
\[
\boxed{
\operatorname{Cov}(X_i,X_j)
=
-\frac{1}{K}h_i h_j^\top,
\qquad i\neq j.
}
\]
More precisely, the scalar covariance of the centred rewards is negative; the
cross-covariance matrix is a negative scalar multiple of $h_i h_j^\top$.
The negative dependence follows from the exact constraint
\[
\sum_{i=1}^K(r_i-\bar r)=0:
\]
if one centred reward increases, the others must collectively decrease.

\paragraph{(iii)}
The full covariance matrix is
\[
\boxed{
\operatorname{Cov}(\hat g)
=
\frac{1}{K^2}
\left[
\sum_{i=1}^K h_i h_i^\top
-
\frac{1}{K}
\left(\sum_{i=1}^K h_i\right)
\left(\sum_{j=1}^K h_j\right)^\top
\right]
=
\frac{1}{K^2}
\sum_{i=1}^K
(h_i-\bar h)(h_i-\bar h)^\top .
}
\]

For the scalar variance convention used in the lectures, define
\[
\operatorname{Var}_{\mathrm{tr}}(Z)
:=
\operatorname{tr}\!\left(\operatorname{Cov}(Z)\right).
\]
Taking the trace gives
\[
\boxed{
\operatorname{Var}_{\mathrm{tr}}(\hat g)
=
\frac{1}{K^2}
\sum_{i=1}^K
\|h_i-\bar h\|^2.
}
\]
Also,
\[
\operatorname{Cov}(X_1)
=
\left(1-\frac{1}{K}\right)h_1h_1^\top,
\]
and hence
\[
\operatorname{Var}_{\mathrm{tr}}(X_1)
=
\left(1-\frac{1}{K}\right)\|h_1\|^2.
\]
If one instead averaged $K$ independent copies of $X_1$, the naive i.i.d.
variance would be
\[
\operatorname{Var}_{\mathrm{tr,iid}}(\hat g)
=
\frac{1}{K}\operatorname{Var}_{\mathrm{tr}}(X_1).
\]

Following the lecture convention, define the effective sample size by
\[
\operatorname{Var}_{\mathrm{tr}}(\hat g)
=
\frac{\operatorname{Var}_{\mathrm{tr}}(X_1)}
{K_{\mathrm{eff}}}.
\]
Therefore
\[
\boxed{
K_{\mathrm{eff}}
=
\frac{
\operatorname{Var}_{\mathrm{tr}}(X_1)
}{
\operatorname{Var}_{\mathrm{tr}}(\hat g)
}
=
\frac{
K^2\left(1-\frac{1}{K}\right)\|h_1\|^2
}{
\sum_{i=1}^K\|h_i-\bar h\|^2
}
=
\frac{
K(K-1)\|h_1\|^2
}{
\sum_{i=1}^K\|h_i-\bar h\|^2
}.
}
\]
Equivalently, its discrepancy from the nominal sample size $K$ is
\[
\frac{K_{\mathrm{eff}}}{K}
=
\frac{
\operatorname{Var}_{\mathrm{tr,iid}}(\hat g)
}{
\operatorname{Var}_{\mathrm{tr}}(\hat g)
}.
\]

If
\[
\frac{1}{K}\sum_{i=1}^K\|h_i-\bar h\|^2
\]
remains of constant order and $\|h_1\|^2=O(1)$, then
\[
K_{\mathrm{eff}}=\Theta(K)
\]
as $K\to\infty$.

If all score vectors lie in one direction, write
\[
h_i=c_i v.
\]
Then
\[
\operatorname{Cov}(\hat g)
=
\frac{vv^\top}{K^2}
\sum_{i=1}^K(c_i-\bar c)^2,
\]
so the covariance has rank at most one, and
\[
K_{\mathrm{eff}}
=
\frac{
K(K-1)c_1^2
}{
\sum_{i=1}^K(c_i-\bar c)^2
}.
\]
Thus only variation in the scalar coefficients $c_i$ contributes to the
variance. In the special case $h_i=h$ for every $i$,
\[
\hat g
=
\frac{h}{K}\sum_{i=1}^K\hat A_i
=0
\]
for every reward realisation. Hence
$\operatorname{Cov}(\hat g)=0$ and, provided $h\neq0$, the effective sample size
is formally infinite.


\subsection*{I.4.2 KL-regularised policy update}

\paragraph{(i)}
The objective is
\[
F(\pi)
=
\sum_a\pi(a)A_t(a)
-\frac{1}{\eta}
\sum_a\pi(a)\log\frac{\pi(a)}{\pi_t(a)}.
\]
Introducing a Lagrange multiplier $\lambda$ for $\sum_a\pi(a)=1$, stationarity
on the support of $\pi_t$ gives
\[
A_t(a)
-\frac{1}{\eta}
\left(
\log\frac{\pi(a)}{\pi_t(a)}+1
\right)
+\lambda=0.
\]
Consequently,
\[
\boxed{
\pi_{t+1}(a)
=
\frac{\pi_t(a)\exp(\eta A_t(a))}
{\sum_b\pi_t(b)\exp(\eta A_t(b))}
}.
\]
Actions outside the support of $\pi_t$ remain assigned zero probability. The
objective is strictly concave on this support, so the optimiser is unique,
assuming the normalising constant is finite.

For small $\eta$,
\[
\pi_{t+1}(a)
=
\pi_t(a)
\left[
1+\eta
\left(
A_t(a)-\mathbb{E}_{\pi_t}[A_t]
\right)
\right]
+O(\eta^2),
\]
so $\eta$ controls the update size. Large $\eta$ places increasingly more mass
on high-advantage actions.

\paragraph{(ii)}
For the true advantage of $\pi_t$,
\[
J(\pi_t)
=
\mathbb{E}_{a\sim\pi_t}[A^{\pi_t}(a)]
=0.
\]
The regularised objective evaluated at $\pi_t$ is also zero. Since
$\pi_{t+1}$ maximises it,
\[
J(\pi_{t+1})
-\frac{1}{\eta}
\operatorname{KL}(\pi_{t+1}\Vert\pi_t)
\geq 0.
\]
It follows that
\[
\boxed{
J(\pi_{t+1})
\geq
\frac{1}{\eta}
\operatorname{KL}(\pi_{t+1}\Vert\pi_t)
\geq 0
=
J(\pi_t).
}
\]
Thus the update guarantees monotone improvement in the scalar performance
measure defined in the question.

In neural-network GRPO the guarantee may fail because the exact advantage is
replaced by a noisy finite-sample estimate and the maximisation is replaced by
a finite number of stochastic-gradient steps in a restricted parametric policy
class. The clipped GRPO objective is also not identical to the exact
KL-regularised objective above.

\paragraph{(iii)}
For one sample, write
\[
\ell(\rho,\hat A)
=
\min\left(
\rho\hat A,
\operatorname{clip}(\rho,1-\epsilon,1+\epsilon)\hat A
\right).
\]
Equivalently,
\[
\ell(\rho,\hat A)
=
\begin{cases}
\hat A\min(\rho,1+\epsilon), & \hat A\geq 0,\\[2mm]
\hat A\max(\rho,1-\epsilon), & \hat A<0.
\end{cases}
\]
Thus, for a positive advantage, the objective gives no further incentive to
increase the ratio above $1+\epsilon$. For a negative advantage, it gives no
further incentive to reduce the ratio below $1-\epsilon$.

This is not a hard constraint
\[
\rho\in[1-\epsilon,1+\epsilon].
\]
Other samples can still move the ratio outside this interval. In contrast, a
constraint
\[
\operatorname{KL}
\bigl(
\pi_\theta(\cdot\mid q)
\Vert
\pi_{\mathrm{old}}(\cdot\mid q)
\bigr)
\leq\delta
\]
is a joint distribution-level constraint coupling all actions. PPO clipping is
samplewise and advantage-dependent, does not directly constrain unsampled
actions, and does not guarantee that the joint KL divergence is small.


\subsection*{I.4.3 Mixture proposal and importance weighting}

Write
\[
p(a)=\pi_{\mathrm{old}}(a\mid q),
\qquad
u(a)=\pi_{\mathrm{unif}}(a\mid q),
\qquad
m(a)=(1-\alpha)p(a)+\alpha u(a),
\]
and let
\[
s_\theta(a)=\nabla_\theta\log\pi_\theta(a\mid q).
\]
For the population analysis, define
\[
V_p=\mathbb{E}_{a\sim p}[R(q,a)],
\qquad
\sigma_p^2=\operatorname{Var}_{a\sim p}[R(q,a)],
\qquad
A_p(a)=\frac{R(q,a)-V_p}{\sigma_p}.
\]

\paragraph{(i)}
The target gradient is
\[
g_p
=
\mathbb{E}_{a\sim p}[s_\theta(a)A_p(a)].
\]
Since $m(a)\geq(1-\alpha)p(a)$, the proposal has support wherever $p$ does.
The importance weight is
\[
\boxed{
w(a)
=
\frac{p(a)}{m(a)}
=
\frac{
\pi_{\mathrm{old}}(a\mid q)
}{
(1-\alpha)\pi_{\mathrm{old}}(a\mid q)
+\alpha\pi_{\mathrm{unif}}(a\mid q)
}.
}
\]
Therefore
\[
g_p
=
\mathbb{E}_{a\sim m}
\left[
w(a)s_\theta(a)A_p(a)
\right],
\]
and an importance-corrected estimator is
\[
\boxed{
\hat g_{\mathrm{mix}}
=
\frac{1}{K}
\sum_{i=1}^K
w_i s_\theta(a_i)\hat A_i^{(p)},
\qquad
w_i=w(a_i),
}
\]
where $\hat A_i^{(p)}$ must estimate the old-policy advantage rather than the
mixture-policy advantage.

\paragraph{(ii)}
As $\alpha\to0$,
\[
m(a)\to p(a),
\qquad
w(a)\to1.
\]
Moreover, the reweighted old-policy moment estimates given below reduce to the
ordinary empirical moments. Hence, for the same realised actions,
\[
\boxed{
\hat g_{\mathrm{mix}}\longrightarrow\hat g_{\mathrm{GRPO}}.
}
\]

\paragraph{(iii-a)}
By linearity of expectation,
\[
\boxed{
V^{\pi_{\mathrm{mix}}}(q)
=
(1-\alpha)V^{\pi_{\mathrm{old}}}(q)
+
\alpha V^{\pi_{\mathrm{unif}}}(q).
}
\]

\paragraph{(iii-b)}
At the population level, for any deterministic baseline $b$,
\[
\mathbb{E}_{m}
\left[
w(a)s_\theta(a)(R(q,a)-b)
\right]
=
\mathbb{E}_{p}
\left[
s_\theta(a)(R(q,a)-b)
\right].
\]
In particular,
\begin{align*}
\mathbb{E}_{m}
\left[
w s_\theta(R-V^{\pi_{\mathrm{mix}}})
\right]
&=
\mathbb{E}_{p}
\left[
s_\theta(R-V^{\pi_{\mathrm{old}}})
\right]\\
&\quad+
\left(
V^{\pi_{\mathrm{old}}}
-
V^{\pi_{\mathrm{mix}}}
\right)
\mathbb{E}_{p}[s_\theta].
\end{align*}
At the policy-gradient evaluation point $\pi_\theta=p$,
\[
\mathbb{E}_{p}[s_\theta]
=
\mathbb{E}_{p}[\nabla_\theta\log p(a)]
=0,
\]
so a deterministic baseline shift alone does not bias the gradient.

However, the empirical $\bar r$ is a random function of the same group and the
empirical $\sigma_r$ is also proposal-dependent. Multiplying only the outer
gradient contribution by $w_i$ therefore does not transform mixture group
normalisation into old-policy group normalisation.

The intended empirical correction is to use weighted moments
\[
S_w=\sum_{j=1}^K w_j,
\qquad
\bar r_w
=
\frac{\sum_{j=1}^K w_jr_j}{S_w},
\]
and
\[
(\sigma_{r,w})^2
=
\frac{
\sum_{j=1}^K w_j(r_j-\bar r_w)^2
}{
S_w
}.
\]
The corresponding reweighted advantage is
\[
\boxed{
\hat A_i^{(w)}
=
\frac{r_i-\bar r_w}{\sigma_{r,w}},
}
\]
giving
\[
\boxed{
\hat g_{\mathrm{mix}}
=
\frac{1}{K}
\sum_{i=1}^K
w_i s_\theta(a_i)
\frac{r_i-\bar r_w}{\sigma_{r,w}}.
}
\]
This self-normalised estimator is consistent for the old-policy moments, but,
like ordinary same-group GRPO normalisation, it is not exactly unbiased at
finite $K$.

For a literally unbiased finite-$K$ baseline correction, one may instead use
the leave-one-out Horvitz--Thompson baseline
\[
\hat V_{p,-i}
=
\frac{1}{K-1}
\sum_{j\neq i}w_jr_j,
\qquad
\mathbb{E}[\hat V_{p,-i}]=V_p.
\]
Since $\hat V_{p,-i}$ is independent of $a_i$,
\[
\mathbb{E}_{m}
\left[
w_i s_\theta(a_i)
\frac{r_i-\hat V_{p,-i}}{\sigma_p}
\right]
=
\mathbb{E}_{p}
\left[
s_\theta(a)
\frac{R(q,a)-V_p}{\sigma_p}
\right].
\]
An estimated standard deviation must similarly be obtained independently if
literal finite-sample unbiasedness is required.

\paragraph{(iv)}
First ignore the additional randomness from estimating the baseline and
standard deviation. Let
\[
Y(a)=s_\theta(a)A_p(a),
\qquad
g_p=\mathbb{E}_{p}[Y(a)].
\]
Then
\[
\operatorname{Var}(\hat g_{\mathrm{GRPO}})
=
\frac{1}{K}
\left(
\mathbb{E}_{p}[YY^\top]-g_pg_p^\top
\right),
\]
whereas
\[
\operatorname{Var}(\hat g_{\mathrm{mix}})
=
\frac{1}{K}
\left(
\mathbb{E}_{m}[w^2YY^\top]-g_pg_p^\top
\right)
=
\frac{1}{K}
\left(
\mathbb{E}_{p}[wYY^\top]-g_pg_p^\top
\right).
\]
Therefore
\[
\boxed{
\operatorname{Var}(\hat g_{\mathrm{mix}})
-
\operatorname{Var}(\hat g_{\mathrm{GRPO}})
=
\frac{1}{K}
\mathbb{E}_{p}
\left[
(w-1)YY^\top
\right].
}
\]
There is no universal ordering between the two covariance matrices. In a
direction $v$,
\[
v^\top
\left[
\operatorname{Var}(\hat g_{\mathrm{mix}})
-
\operatorname{Var}(\hat g_{\mathrm{GRPO}})
\right]v
=
\frac{1}{K}
\mathbb{E}_{p}
\left[
(w-1)(v^\top Y)^2
\right].
\]
Moreover,
\[
w(a)-1
=
\frac{\alpha(p(a)-u(a))}{m(a)}.
\]
Thus mixing increases variance when large squared gradient contributions,
and hence large squared reward deviations, are concentrated on responses for
which $p(a)>u(a)$. Those responses are undersampled by the mixture relative to
$p$ and receive weights $w(a)>1$.

Conversely, mixing decreases variance when large reward deviations are
concentrated on responses for which $u(a)>p(a)$. The uniform component
oversamples these responses and their weights satisfy $w(a)<1$.

The change in the unweighted reward variance can also be seen directly:
\[
\operatorname{Var}_{m}(R)
=
(1-\alpha)\operatorname{Var}_{p}(R)
+
\alpha\operatorname{Var}_{u}(R)
+
\alpha(1-\alpha)(V_p-V_u)^2.
\]
Hence mixture sampling increases the raw reward variance if
\[
\operatorname{Var}_{u}(R)
+
(1-\alpha)(V_p-V_u)^2
>
\operatorname{Var}_{p}(R),
\]
and decreases it if the reverse inequality holds. Importance weighting adds
the separate $w^2$ effect described above.

\paragraph{(v-a)}
Under the stated assumption,
\[
w_i=c^T=\exp(T\log c).
\]
Since $0<c<1$, $\log c<0$, and therefore
\[
\boxed{
w_i\to0
\quad\text{exponentially fast as }T\to\infty.
}
\]
A ratio below one is typical for exploratory tokens or responses that are
unlikely under $\pi_{\mathrm{old}}$ but receive additional probability from
the uniform component, so that
$\pi_{\mathrm{mix}}>\pi_{\mathrm{old}}$. It is not implied for every token
merely by $\alpha>0$.

\paragraph{(v-b)}
Such trajectories make exponentially small contributions to the estimator.
If many sampled trajectories have collapsed weights, the estimate is
effectively based on only a small number of old-policy-like trajectories,
causing a small effective sample size and unstable estimates; the products may
also underflow numerically.

One unbiased remedy is to keep the proposal closer to the old policy, for
example by reducing $\alpha$ or enforcing the per-token log-ratio condition
\[
\left|
\log
\frac{
\pi_{\mathrm{old}}(a_t\mid q,a_{<t})
}{
\pi_{\mathrm{mix}}(a_t\mid q,a_{<t})
}
\right|
\leq\frac{\delta}{T}.
\]
Then
\[
e^{-\delta}\leq w_i\leq e^\delta,
\]
so the full-response weight cannot collapse exponentially. The trade-off is
weaker uniform exploration.

Alternatively, clipping or self-normalising the response weights reduces
variance and numerical collapse, but introduces bias.